\documentclass[11pt,a4paper]{article}
\usepackage[T1]{fontenc}
\usepackage{lmodern,amsmath,amssymb,amsthm,mathtools,booktabs,microtype}
\usepackage[margin=27mm]{geometry}
\usepackage{enumitem}
\usepackage{xcolor}
\definecolor{ink}{HTML}{324D54}
\usepackage[colorlinks=true,linkcolor=ink,urlcolor=ink,citecolor=ink]{hyperref}
\usepackage[nameinlink,noabbrev,capitalise]{cleveref}
\setlength{\emergencystretch}{1.5em}
\numberwithin{equation}{section}
\newtheorem{theorem}{Theorem}[section]
\newtheorem{lemma}[theorem]{Lemma}
\newtheorem{corollary}[theorem]{Corollary}
\newtheorem{proposition}[theorem]{Proposition}
\newtheorem{conjecture}[theorem]{Conjecture}
\theoremstyle{remark}
\newtheorem{remark}[theorem]{Remark}
\crefname{conjecture}{Conjecture}{Conjectures}
\newcommand{\E}{\mathbb E}
\newcommand{\PP}{\mathbb P}
\newcommand{\R}{\mathbb R}
\newcommand{\cH}{\mathcal H}
\newcommand{\dist}{\operatorname{dist}}
\newcommand{\Var}{\operatorname{Var}}
\newcommand{\Ent}{\operatorname{Ent}}
\title{A sharp stop-loss bound for Rademacher sums,\\ the constant two in a concave comparison,\\ and local extremal sets on the discrete cube}
\author{Anonymous\\[3pt]{\small DOI \href{https://doi.org/10.5281/zenodo.24548728}{10.5281/zenodo.24548728}}}
\date{25 November 2026}
\begin{document}
\maketitle

\begin{abstract}
Flip $n$ fair coins, attach real weights with squares summing to one, and
add the signed weights. How much of this random sum can lie more than one
standard deviation from zero, on average? We prove that the mean excess
$\E(|S|-1)_+$ is at most $(\sqrt2-1)/2$, with equality only when two weights
are equal and all the others vanish, and with a quadratic stability estimate.
This settles a bound proposed by Lowther on MathOverflow and, through his
reduction, Kontorovich's question: for independent symmetric random variables
and every increasing concave $f$ with $f(0)=0$, the expectation of $f$ at
the square function is at most twice its expectation at the absolute sum.
The constant two is sharp, the comparison extends to sign-invariant vectors in
any Hilbert space with linear stability, and it improves the universal constant
in Kontorovich's lower bound for the total variation distance between product
measures from $0.1798$ to $(1-e^{-1/2})/2\approx0.1967$.
The same inequality is the support function of the set of centroids of
subsets of the discrete cube. Through it we study sets of density near
$1/4$ with maximal first-level Fourier weight, the equivalent of the
Ahlswede--Katona minimum average distance problem. Without further
input we show that the dimension-free maximal weight $W$ satisfies
$W(1/4\pm h)=2\min(1/4,1/4\pm h)^2+(2+o(1))h^2\log(1/h)$. Assuming
Lei Yu's recent small-density theorem (an unrefereed 2026 preprint), we
obtain an exact finite-parameter formula near $1/4$, the next-order terms,
and nonattainment in finite dimension. Above $1/4$ the optimal sets split
added mass between two neighbouring faces in proportions $0.979:0.021$,
the solution of a Curie--Weiss--Potts variational problem. An extension of
Fang and Huang's capped three-coefficient inequality to all coefficients,
with stability, gives the analogous results near density $1/8$, again
assuming Yu's theorem.
\end{abstract}

\tableofcontents

%======================================================================
\section{Introduction}\label{sec:intro}

This section explains the question, why it matters, what was known before,
what this paper adds, and where each result is proved.

\paragraph{The question.}
Let $\varepsilon_1,\ldots,\varepsilon_n$ be independent random signs, each
equal to $+1$ or $-1$ with probability one half, and let $a\in\R^n$ be a
vector of real weights with $\sum_ia_i^2=1$. The weighted sum
\[
S_a=\sum_{i=1}^n a_i\varepsilon_i
\]
is called a Rademacher sum; it has mean zero and variance one. In September
2024 Kontorovich asked on MathOverflow \cite{Kontorovich} whether
\begin{equation}\label{eq:question}
 \E f\Bigl(\sqrt{\textstyle\sum_i X_i^2}\Bigr)\le 2\,\E f\Bigl(\Bigl|\textstyle\sum_i X_i\Bigr|\Bigr)
\end{equation}
for all independent symmetric real random variables $X_1,\ldots,X_n$ and every
increasing concave function $f$ with $f(0)=0$. The left side measures the
size of the summands through their square function; the right side measures
the size of their sum. Lowther \cite{Lowther} reduced the question to Rademacher
sums of unit variance and to the single function $f(t)=\min(t,1)$, so that
\eqref{eq:question} is equivalent to
\begin{equation}\label{eq:clip}
 \E\min(|S_a|,1)\ge\tfrac12\qquad\text{for every unit vector }a .
\end{equation}
He proved \eqref{eq:question} with the constant $(1/\sqrt2-1/4)^{-1}\approx2.187$
and observed that the constant $2$ would follow from the mean-excess bound
\begin{equation}\label{eq:lowther}
 \E(|S_a|-1)_+\le\frac{\sqrt2-1}{2},
\end{equation}
where $x_+=\max(x,0)$. Both bounds are attained when two weights equal
$1/\sqrt2$ and the others vanish; then $|S_a|$ is $0$ or $\sqrt2$ with equal
probabilities. The thread records the constant two as open.

\paragraph{Why it matters.}
Kontorovich's motivation was the total variation distance between product
measures. His lower bound for that distance in terms of the marginal
distances \cite{KontorovichTV} uses \eqref{eq:question} with Lowther's
constant; a sharp constant improves it directly
(\cref{cor:tv}). More broadly, $t\mapsto\E(S_a-t)_+$ is the
\emph{stop-loss transform} of $S_a$, a standard object in the comparison
of random variables (see \cite[Section~II]{YuLevel}). Sharp lower bounds for
$\E|S_a|$ (Khintchine's inequality, with Szarek's constant $1/\sqrt2$
\cite{Szarek}) and for the tail probabilities $\PP(|S_a|\ge t)$
\cite{HP} are classical or recent landmarks. A sharp bound for the
stop-loss transform near one standard deviation, where two equal weights are
extremal, is a natural companion to these results.

The same functional also governs a problem on the discrete cube
$\{-1,1\}^n$. For a set $A$ of density $\mu(A)$, write
$b_A=\E[\varepsilon\,1_A]\in\R^n$ for its unnormalised centroid. Its
squared length $\|b_A\|^2$ is the first-level Fourier weight of the
indicator $1_A$, and maximising it at fixed density is equivalent to
minimising the average Hamming distance within $A$, a problem posed by
Ahlswede and Katona \cite{AhlswedeKatona} (see \cite[Section~1]{FangHuang} and \cite[Section~I]{YuLevel}). Since
$\|b_A\|=\max_{\|a\|=1}\E[S_a1_A]$, every bound on $\E(S_a-t)_+$ bounds
$\|b_A\|-t\mu(A)$; this is how Fang and Huang settled density $1/8$
\cite{FangHuang}.

\paragraph{What was known.}
Szarek proved $\E|S_a|\ge1/\sqrt2$ \cite{Szarek}, and Haagerup found all
the sharp Khintchine constants \cite{Haagerup}. De, Diakonikolas and
Servedio \cite{DDS} proved a robust version: $\E|S_a|-1/\sqrt2$ is at least
a constant multiple of the distance from $a$ to the two-equal-weight vectors;
Eskenazis, Nayar and Tkocz made the constant explicit \cite{ENT}.
Hollom and Portier determined $\inf_a\PP(|S_a|\ge y)$ for every $y$,
confirming another conjecture of Lowther; the optimal weights change with
$y$ \cite{HP}. Yu proved that, for thresholds $t$ beyond a large universal
constant, $\E(S_a-t)_+$ is strictly below its Gaussian value
\cite[Theorem~4]{YuLevel}.
On the cube, the maximal first-level weight at densities $1/2$ and $1/4$ is
attained by subcubes \cite{AlthoferSillke,FuWeiYeung} (see the account in \cite[Section~1]{FangHuang}); Fang and Huang proved the
same at density $1/8$ \cite{FangHuang}; and Yu proved that at all
sufficiently small densities the dimension-free maximum equals the Gaussian
(Hamming-ball) value and is not attained in any finite dimension
\cite[Theorem~5]{YuLevel}. Yu asked for the value at $1/8$
\cite[Problem~2]{YuLevel}; Fang and Huang conjectured the inequality that
would settle $1/16$ \cite[Conjecture~3.1]{FangHuang}.

\paragraph{What is new.}
The main results are as follows. Items 1--6 are unconditional.
Items 7 and 8 use Yu's Theorem~5 \cite{YuLevel}, a September 2026 preprint
that has not yet been refereed.
\begin{enumerate}[leftmargin=2em]
\item \textbf{Sharp stop-loss bound} (\cref{thm:stoploss}). For every unit
 vector $a$ and every $t\in[1/\sqrt2,\,201/200]$,
 $\E(|S_a|-t)_+\le(\sqrt2-t)/2$, with equality only for two equal weights.
 At $t=1$ this is Lowther's bound \eqref{eq:lowther}, with a deficit of at
 least $\dist(a,\mathcal E_2)^2/1000$. The interval cannot extend beyond
 $[2-\sqrt2,\,2\sqrt2-\sqrt3]$ (\cref{prop:interval}), and numerical
 evidence suggests that this larger interval is the true one
 (\cref{conj:interval}).
\item \textbf{The constant two} (\cref{cor:clip,thm:hilbert}). Inequality
 \eqref{eq:clip} holds, so \eqref{eq:question} holds with the sharp
 constant $2$, for every nondecreasing $f$ with $f(0)=0$ and $f(t)/t$
 nonincreasing, and for sign-invariant vectors in any Hilbert space. As a
 by-product, $\inf_a\E\min(|S_a|,t)=t/2$ for $0<t\le201/200$
 (\cref{cor:truncated}).
\item \textbf{Product measures} (\cref{cor:tv}). The constant in Kontorovich's
 lower bound $\|P^{\otimes}-Q^{\otimes}\|_{TV}\ge c\min\{1,\|\delta\|_2\}$
 improves from $0.1798$ to $(1-e^{-1/2})/2$.
\item \textbf{Stability} (\cref{thm:stability}). In any Hilbert space,
 $\E\min(\|\sum\varepsilon_ix_i\|,1)-\tfrac12$ is comparable, up to
 absolute constants, to the distance from the tuple $(x_i)$ to two equal
 collinear vectors. The constants are explicit (\cref{rem:explicit}).
\item \textbf{Three equal weights} (\cref{thm:triadic}). Fang and Huang's
 inequality $\E(|S|-2)_+\le1/4$ for variance-three sums with weights at most
 one \cite[Lemma~2.4]{FangHuang} holds without the cap on the weights, with
 a quadratic stability estimate and in a neighbourhood of the threshold.
\item \textbf{Local profile near density $1/4$, leading term}
 (\cref{prop:cubestable,thm:cusp}). Sets of nearly maximal first-level weight
 near density $1/4$ are close to two-coordinate faces, and
 $W(1/4\pm h)=2\min(1/4,1/4\pm h)^2+(2+o(1))h^2\log(1/h)$.
\item \textbf{Local profile, finer structure} (\cref{cor:fourfibres,thm:refined}).
 Using Yu's theorem: an exact formula for $W$ near $1/4$ with three real
 parameters, nonattainment in finite dimension, and the expansion through
 the $h^2$ term. Above $1/4$ the added mass is split between two
 neighbouring faces in the ratio $\lambda_*:(1-\lambda_*)$ with
 $\lambda_*=0.97875\ldots$.
\item \textbf{A transfer principle} (\cref{thm:faces}). Any stable stop-loss
 inequality with $k$ equal extremal weights yields, using Yu's theorem,
 the same local theory near density $2^{-k}$. It applies for $k=2$ and, by
 item~5, for $k=3$ (\cref{cor:eighth}). Numerically it cannot apply for
 $k=5,6$ (\cref{rem:scope}).
\end{enumerate}
The proofs of items 1 and 5 are short case analyses built on a quadratic
majorant with a cosine correction, an idea also used by Fang and Huang
\cite{FangHuang}. The cube results combine the stability estimates with a
Gaussian construction in the remaining coordinates.

\paragraph{Roadmap.}
\Cref{sec:notation} fixes notation. \Cref{sec:excess} proves the sharp
stop-loss bound. \Cref{sec:comparison} derives the clipped bound, the
concave comparison in Hilbert space and the application to product
measures. \Cref{sec:stability} proves stability. \Cref{sec:triadic} treats
three equal weights. \Cref{sec:cube} studies the cube near density $1/4$,
first without and then with Yu's theorem, and \cref{sec:faces} gives the
general transfer theorem and density $1/8$. \Cref{sec:computations}
describes the accompanying computations, and \cref{sec:open} lists open
problems.

%======================================================================
\section{Notation and conventions}\label{sec:notation}

This section collects the notation used throughout.

The signs $\varepsilon=(\varepsilon_1,\ldots,\varepsilon_n)$ are independent
and uniform on $\{-1,1\}$; $\E$ and $\PP$ refer to them unless stated
otherwise. For $a\in\R^n$ we write $S_a=\sum_ia_i\varepsilon_i$ and
$\|a\|$ for the Euclidean norm. A \emph{unit vector} has $\|a\|=1$.
Zero coordinates may always be appended, so we assume $n\ge k$ when
comparing with vectors that have $k$ nonzero entries.

For $k\ge1$, $\mathcal E_k$ denotes the set of \emph{signed permutations} of
$(1,\ldots,1,0,\ldots,0)/\sqrt k$ with exactly $k$ nonzero entries (in the
current dimension). For a unit vector $a$ we write
$d(a)=\dist(a,\mathcal E_2)$ for its Euclidean distance to the
two-equal-weight vectors. If $|a_1|\ge|a_2|$ are the two largest
magnitudes, then $d(a)^2=2-\sqrt2(|a_1|+|a_2|)$.

$\varphi$ and $\Phi$ are the standard normal density and distribution
function, $G$ is a standard normal variable, and
$I(p)=\varphi(\Phi^{-1}(p))$ for $0<p<1$, with $I(0)=I(1)=0$. The function
$I$ is the Gaussian isoperimetric profile: if $T$ is standard normal and
$\PP(T>\tau)=p$, then $\E[T1_{\{T>\tau\}}]=I(p)$.

$\cH$ is a real Hilbert space with norm $\|\cdot\|$; complex spaces are
covered by using the real part of the inner product. Entropy is
$\Ent(\lambda)=-\sum_i\lambda_i\log\lambda_i$ for a probability vector
$\lambda$ (natural logarithms, $0\log0=0$), and $\mathcal P_k$ is the set of
probability vectors of length $k$.

On the cube $\{-1,1\}^n$ with uniform measure $\mu$, a set $A$ has
centroid vector $b_A=\E[\varepsilon1_A]$. A \emph{$k$-face} is a set
$\{\varepsilon_{i_1}=\sigma_1,\ldots,\varepsilon_{i_k}=\sigma_k\}$ obtained by
fixing $k$ coordinates; it has density $2^{-k}$.

%======================================================================
\section{The sharp stop-loss bound}\label{sec:excess}

This section proves the main inequality, \eqref{eq:lowther} and its
extension to an interval of thresholds, and explains why the interval cannot
be much larger.

\begin{theorem}\label{thm:stoploss}
For every $t$ with $1/\sqrt2\le t\le 201/200$,
\[
 \sup_{n,\ \|a\|=1}\E(|S_a|-t)_+=\frac{\sqrt2-t}{2}.
\]
Equality holds exactly when $a\in\mathcal E_2$, that is, when two weights
have magnitude $1/\sqrt2$ and the others vanish. At $t=1$,
\begin{equation}\label{eq:stab1}
 \frac{\sqrt2-1}{2}-\E(|S_a|-1)_+\ \ge\ \frac{d(a)^2}{1000},
\end{equation}
and the quadratic order is best possible.
\end{theorem}

The value at $t=1$ is Lowther's proposed bound \eqref{eq:lowther}. The proof
has three steps: a cosine bound that handles all weight vectors except those
with two large weights; a direct calculation for two large weights; and a
convexity argument for thresholds below one.

\begin{proof}
\emph{Step 1: a cosine bound.} For real $x$,
\[
 \frac{x^2}{4}-(|x|-1)_+=\frac{\dist(x,\{0,2,-2\})^2}{4}
 \ge\frac{1-\cos(\pi x)}{2\pi^2}.
\]
The identity is checked on $|x|\le1$ and $|x|\ge1$ separately; the
inequality follows from $1-\cos y\le y^2/2$, because the three points are
even integers and $\cos(\pi x)$ has period two. Put
$P=\E\cos(\pi S_a)=\prod_i\cos(\pi a_i)$. Taking expectations and using
$\E S_a^2=1$,
\begin{equation}\label{eq:cos-stop}
 \E(|S_a|-1)_+\le\frac14-\frac{1-P}{2\pi^2}.
\end{equation}

\emph{Step 2: most weight vectors.} Order the magnitudes decreasingly and
let $s=|a_1|+|a_2|$, $u=|a_1|-|a_2|$. Suppose first that we are \emph{not}
in the case ``exactly two magnitudes exceed $1/2$ and $s\ge6/5$''. If no
magnitude exceeds $1/2$, then $|P|\le\exp(-\pi^2/2)<1/10$, by
$|\cos z|\le e^{-z^2/2}$ for $|z|\le\pi/2$. Magnitudes in $(1/2,1]$ have
negative cosines and the others nonnegative cosines, so an odd number of
them gives $P\le0$; four of them are impossible since their squares would
sum to more than one. If exactly two exceed $1/2$ and $s<6/5$, the product
of their cosines is
$\tfrac12[\cos(\pi s)+\cos(\pi u)]\le\cos^2(\pi s/2)=\sin^2(\pi(s-1)/2)
\le\sin^2(\pi/10)<1/10$. In all these cases $P\le1/10$, and
\eqref{eq:cos-stop} gives
\[
 \frac{\sqrt2-1}{2}-\E(|S_a|-1)_+\ \ge\ \frac{\sqrt2-1}{2}-\frac14+\frac{9}{20\pi^2}
 =0.0027013\ldots>\frac1{400},
\]
using $\sqrt2>1.4142$ and $\pi^2<9.88$. The derivative in $t$ of
$(\sqrt2-t)/2-\E(|S_a|-t)_+$ is $-\tfrac12+\PP(|S_a|>t)$, which lies in
$[-\tfrac12,\tfrac12]$. Hence at every $t\in[1,201/200]$ the deficit
remains at least $0.0027013-0.0025>1.4\cdot10^{-4}$, and on
$[1/\sqrt2,1]$ it is handled in Step~4. At $t=1$, since $d(a)^2\le2$, the
deficit exceeds $1/400>d(a)^2/1000$.

\emph{Step 3: two large weights.} It remains to treat $|a_1|,|a_2|>1/2$ with
$s\ge6/5$; write $a=|a_1|$, $b=|a_2|$ and let $R$ be the sum of the
remaining terms, with variance $v=1-a^2-b^2=1-(s^2+u^2)/2$. Let
$1\le t\le201/200$; then $0\le u<t<s$. For a constant $c$ and the symmetric
variable $R$ one has $\E(c+R)_+=c_++\tfrac12\E(|R|-|c|)_+$. Applying this to
the four equally likely values $\pm s,\pm u$ of $|a_1|\varepsilon_1'+|a_2|\varepsilon_2'$ (after sign changes) gives
\begin{equation}\label{eq:four-hinges}
 \E(|S_a|-t)_+=\frac{s-t}{2}+
 \frac14\sum_{k\in\{s-t,\ t-u,\ t+u,\ s+t\}}\E(|R|-k)_+ .
\end{equation}
The four cutoffs are positive and already increasing, because
$s+u=2a\le2\le2t$. For increasing positive $k_1\le\cdots\le k_4$ with
partial sums $K_j=k_1+\cdots+k_j$, the sum of the hinges $(x-k_i)_+$ at
$x\ge0$ equals $\max_{0\le j\le4}(jx-K_j)$, and $jx-K_j\le j^2x^2/(4K_j)$ by
the arithmetic--geometric mean inequality. Here
$K=(s-t,\ s-u,\ s+t,\ 2(s+t))$, so the hinge sum is at most $Cx^2$ with
$C=\max_jj^2/(4K_j)$. Using $s-t\ge6/5-201/200=39/200$, $s-u=2b>1$ and
$s+t\ge11/5$ we get $C\le50/39$, and $C\le5/4$ when $t=1$. Since
$v\le1-s^2/2=(\sqrt2-s)(\sqrt2+s)/2\le\sqrt2(\sqrt2-s)$,
\[
 \frac{\sqrt2-t}{2}-\E(|S_a|-t)_+\ \ge\ \frac{\sqrt2-s}{2}-\frac{Cv}{4}
 \ \ge\ \Bigl(\frac12-\frac{\sqrt2\,C}{4}\Bigr)(\sqrt2-s).
\]
The bracket is at least $1/2-50\sqrt2/156=0.0467\ldots>0$ on the interval,
and $1/2-5\sqrt2/16=0.0580\ldots$ at $t=1$. Since
$\sqrt2-s=d(a)^2/\sqrt2$, the deficit at $t=1$ is at least
$0.041\,d(a)^2$, which gives \eqref{eq:stab1}. The deficit vanishes only if
$s=\sqrt2$, which forces $a=b=1/\sqrt2$ and $v=0$.
For sharpness, take $a_1=\cos\theta$, $a_2=\sin\theta$ near
$\theta=\pi/4$ and no other weights: then
$\E(|S_a|-1)_+=(s-1)/2$, so the deficit is
$(\sqrt2-s)/2=d(a)^2/(2\sqrt2)$.

\emph{Step 4: thresholds below one.} At $t=1/\sqrt2$, the bound
$(x-t)_+\le x^2/(4t)$ gives $\E(|S_a|-t)_+\le1/(4t)=(\sqrt2-t)/2$. The map
$t\mapsto\E(|S_a|-t)_+$ is convex, so on $[1/\sqrt2,1]$ it lies below the
chord through its values at the endpoints, which is the line
$(\sqrt2-t)/2$; the pair attains this line. Equality at $t=1/\sqrt2$ forces
$|S_a|\in\{0,\sqrt2\}$ almost surely. If $a$ has $k$ nonzero weights, flip
their signs one by one from all minus to all plus (after making all weights
positive): this produces $k+1$ strictly increasing values of $S_a$, so
$k\le2$, and then $|S_a|\in\{0,\sqrt2\}$ forces two equal weights. At an
interior point of the chord, equality forces equality at $t=1$, which was
characterised in Steps 2--3.
\end{proof}

\begin{remark}
The bound \eqref{eq:cos-stop} is a quadratic majorant with a trigonometric
correction whose expectation factorises over the signs. Fang and Huang use
a majorant of the same type, with two harmonics, for three equal weights
\cite[Lemma~2.3]{FangHuang}; we use theirs in \cref{sec:triadic}.
\end{remark}

The interval $[1/\sqrt2,201/200]$ is where our proof works. The next
proposition shows that the pair cannot be optimal on a much larger interval.

\begin{proposition}\label[proposition]{prop:interval}
If $0\le t<2-\sqrt2$, a single weight ($a=e_1$) gives
$\E(|S_a|-t)_+=1-t>(\sqrt2-t)/2$. If $2\sqrt2-\sqrt3<t<\sqrt3$, three equal
weights give $\E(|S_a|-t)_+=(\sqrt3-t)/4>(\sqrt2-t)/2$. Hence the largest
interval of thresholds on which two equal weights maximise the stop-loss
transform is contained in $[2-\sqrt2,\,2\sqrt2-\sqrt3]\approx[0.586,1.096]$.
\end{proposition}

\begin{proof}
For $a=e_1$, $|S_a|=1$. For three weights $1/\sqrt3$, $|S_a|$ equals
$1/\sqrt3$ with probability $3/4$ and $\sqrt3$ with probability $1/4$. The
stated inequalities are linear in $t$.
\end{proof}

We conjecture that this interval is exact; see \cref{conj:interval} and
the evidence in \cref{sec:computations}. The proof above does not reach the
endpoints because two different weight vectors tie there, so the margin
used in Step~2 disappears.

%======================================================================
\section{The constant two and its consequences}\label{sec:comparison}

This section derives the clipped bound \eqref{eq:clip}, the concave
comparison \eqref{eq:question} with constant two in any Hilbert space, a
sharp truncated-mean inequality, and the improvement in Kontorovich's
product-measure bound.

\begin{corollary}[Clipped first moment]\label[corollary]{cor:clip}
For every unit vector $a$, $\E\min(|S_a|,1)\ge1/2$, with equality exactly
when $a\in\mathcal E_2$.
\end{corollary}

\begin{proof}
Since $\min(x,1)=x-(x-1)_+$ for $x\ge0$, Szarek's inequality
$\E|S_a|\ge1/\sqrt2$ \cite{Szarek} (see also \cite[Theorem~1.1]{Haagerup}) and
\cref{thm:stoploss} at $t=1$ give
$\E\min(|S_a|,1)\ge1/\sqrt2-(\sqrt2-1)/2=1/2$. Equality requires equality in
\cref{thm:stoploss}.
\end{proof}

\begin{remark}[A second proof]\label[remark]{rem:second}
\Cref{cor:clip} also has a short direct proof. Let $m=\max_i|a_i|$. If
$m\ge1/2$, condition on all signs except that of a largest weight; the two
resulting values of $S_a$ differ by $2m\ge1$, and the triangle inequality
for the bounded metric $\min(|x-y|,1)$ gives
$\min(|x+m|,1)+\min(|x-m|,1)\ge1$. If $m\le1/2$, Haagerup's coefficient
inequality $\E|S_a|\ge\sum_ia_i^2F(a_i^{-2})$, with
$F(s)=2\Gamma((s+1)/2)/(\sqrt{\pi s}\,\Gamma(s/2))$ increasing
\cite[Lemmas~1.3 and~1.4]{Haagerup} (see also
\cite[Section~4, equation~(9)]{Melbourne} and \cite[equation~(62)]{ENT}),
gives $\E|S_a|\ge F(4)=3/4$, and $(x-1)_+\le x^2/4$ gives
$\E\min(|S_a|,1)\ge3/4-1/4$. The equality analysis is a short additional
argument.
\end{remark}

\begin{corollary}[Sharp truncated mean]\label[corollary]{cor:truncated}
For $0<t\le201/200$,
\[
 \inf_{n,\ \|a\|=1}\E\min(|S_a|,t)=\frac t2,
\]
attained exactly by $a\in\mathcal E_2$.
\end{corollary}

\begin{proof}
For $t\le1$, the pointwise bound $\min(y,t)\ge t\min(y,1)$ and
\cref{cor:clip} give $\E\min(|S_a|,t)\ge t/2$. For $1<t\le201/200$, Szarek's
inequality and \cref{thm:stoploss} give
$\E\min(|S_a|,t)=\E|S_a|-\E(|S_a|-t)_+\ge1/\sqrt2-(\sqrt2-t)/2=t/2$. For
$a\in\mathcal E_2$, $\E\min(|S_a|,t)=\tfrac12\min(\sqrt2,t)=t/2$.
\end{proof}

A sharp bound for each tail probability $\PP(|S_a|\ge y)$ would not give
\cref{cor:truncated} by integration over $y$, because the optimal weights
change with $y$ \cite[Theorem~1.1]{HP}.

We now pass to vectors. Call the joint law of $(X_1,\ldots,X_n)$ in $\cH$
\emph{sign invariant} if multiplying the coordinates by any fixed signs
preserves the law. Independent symmetric random vectors have this property.

\begin{theorem}[Concave comparison]\label{thm:hilbert}
Let $(X_1,\ldots,X_n)$ be a sign-invariant tuple in a Hilbert space $\cH$.
Let $f\colon[0,\infty)\to[0,\infty)$ be nondecreasing with $f(0)=0$ and
$f(t)/t$ nonincreasing on $(0,\infty)$. Then
\[
 \E f\Bigl(\sqrt{\textstyle\sum_i\|X_i\|^2}\Bigr)
 \le2\,\E f\Bigl(\Bigl\|\textstyle\sum_iX_i\Bigr\|\Bigr),
\]
where both sides may be $+\infty$. Every nondecreasing concave $f$ with
$f(0)=0$ qualifies. The constant $2$ is sharp.
\end{theorem}

\begin{proof}
\emph{Deterministic vectors.} Fix $x_1,\ldots,x_n\in\cH$ and put
$s^2=\sum_i\|x_i\|^2>0$. The normalised Gram matrix
$A=(\langle x_i,x_j\rangle/s^2)$ is positive semidefinite with trace one,
so $A=\sum_kp_kv_kv_k^{\mathsf T}$ with unit vectors $v_k\in\R^n$,
$p_k>0$ and $\sum_kp_k=1$. Then
$\|\sum_i\varepsilon_ix_i\|^2/s^2=\varepsilon^{\mathsf T}A\varepsilon
=\sum_kp_k(v_k\cdot\varepsilon)^2$ and
\begin{equation}\label{eq:lift}
 \min\Bigl(\sqrt{\varepsilon^{\mathsf T}A\varepsilon},1\Bigr)
 \ge\Bigl(\sum_kp_k\min((v_k\cdot\varepsilon)^2,1)\Bigr)^{1/2}
 \ge\sum_kp_k\min(|v_k\cdot\varepsilon|,1).
\end{equation}
The first inequality holds because $\sum_kp_k\min(b_k,1)\le\min(\sum_kp_kb_k,1)$;
the second is Jensen's inequality. Taking expectations and applying
\cref{cor:clip} to each $v_k$ gives
$\E\min(\|\sum_i\varepsilon_ix_i\|/s,1)\ge1/2$. The hypotheses on $f$ give
$f(y)\ge f(s)\min(y/s,1)$ for all $y\ge0$ (use monotonicity of $f$ for
$y\ge s$ and of $f(t)/t$ for $y<s$). Therefore
$\E f(\|\sum_i\varepsilon_ix_i\|)\ge\tfrac12f(s)$.

\emph{Random vectors.} Let $\varepsilon$ be independent of $X$. Apply the
deterministic inequality conditionally on $X$ and integrate. Sign invariance
shows that $\sum_i\varepsilon_iX_i$ has the same law as $\sum_iX_i$, while
the square function is unchanged.

\emph{Sharpness.} For $X_1,X_2$ independent and uniform on
$\{\pm1/\sqrt2\}$ and $f(t)=\min(t,1)$, the left side is $1$ and the right
side is $2\cdot\tfrac12$.
\end{proof}

\Cref{thm:hilbert} answers Kontorovich's question \eqref{eq:question}. The
class of $f$ is larger than the concave class: for example
$f(t)=\min(t,1)+\tfrac12\min((t-2)_+,2)$ is not concave but qualifies:
it is nondecreasing, and $f(t)/t$ equals $1$, $1/t$, $1/2$ and $2/t$ on
$[0,1]$, $[1,2]$, $[2,4]$ and $[4,\infty)$.

\begin{corollary}[Product measures]\label[corollary]{cor:tv}
Let $P_1,\ldots,P_n$ and $Q_1,\ldots,Q_n$ be probability measures on finite
sets, and let $\delta_i=\|P_i-Q_i\|_{TV}$. Then
\[
 \Bigl\|\bigotimes_iP_i-\bigotimes_iQ_i\Bigr\|_{TV}
 \ \ge\ \frac{1-e^{-1/2}}{2}\,\min\{1,\|\delta\|_2\},
 \qquad \frac{1-e^{-1/2}}{2}=0.19673\ldots
\]
\end{corollary}

\begin{proof}
We follow Kontorovich's proof of \cite[Theorem~1]{KontorovichTV} and change
one step. By his Lemmas~2 and~3 (reduction to Bernoulli measures and
symmetrisation) it suffices to treat $\hat P=\mathrm{Ber}(\hat p)$ and
$\hat Q=\mathrm{Ber}(\hat q)$ on $\{0,1\}^n$ with $\hat q=1-\hat p$ and
$\|\hat p-\hat q\|_2\ge\tfrac12\|\delta\|_2$; the distance can only
decrease in this reduction. With $Z$ uniform on $\{0,1\}^n$ and
$L_i=\log(\hat P_i(Z_i)/\hat Q_i(Z_i))$, his computation gives
\[
 \|\hat P-\hat Q\|_{TV}\ \ge\ \E\Bigl[1-\exp\Bigl(-\tfrac12\Bigl|\sum_iL_i\Bigr|\Bigr)\Bigr].
\]
The $L_i$ are independent and symmetric (because $\hat q=1-\hat p$), and
$|L_i|=\log(\hat p_i/\hat q_i)$ in absolute value is constant. Apply
\cref{thm:hilbert} with $\cH=\R$ and the concave function
$f(t)=1-e^{-t/2}$; where Kontorovich used Lowther's factor
$1/\sqrt2-1/4$, we obtain the factor $1/2$:
\[
 \|\hat P-\hat Q\|_{TV}\ge\tfrac12 f\Bigl(\sqrt{\textstyle\sum_iL_i^2}\Bigr)
 \ge\tfrac12\bigl(1-e^{-\|\hat p-\hat q\|_2}\bigr),
\]
using $|\log((1+x)/(1-x))|\ge2|x|$ with $x=\hat p_i-\hat q_i$. Finally
$1-e^{-y}\ge(2-2e^{-1/2})\min(y,\tfrac12)$ for $y\ge0$, by concavity, and
$\min(\|\hat p-\hat q\|_2,\tfrac12)\ge\tfrac12\min(\|\delta\|_2,1)$.
\end{proof}

Kontorovich's printed constant is $c>0.1798$, which is
$(1/\sqrt2-1/4)(1-e^{-1/2})=0.17985\ldots$ \cite[Theorem~1]{KontorovichTV}.

%======================================================================
\section{Stability}\label{sec:stability}

This section shows that near-equality in \cref{thm:hilbert} forces the
vectors to be close to two equal collinear vectors, with a linear rate.

For a tuple $x=(x_1,\ldots,x_n)$ in $\cH$ with $\sum_i\|x_i\|^2=1$ put
\[
 D(x)=\E\min\Bigl(\Bigl\|\sum_i\varepsilon_ix_i\Bigr\|,1\Bigr)-\frac12,
 \qquad
 d_\cH(x)=\min_{w\in\mathcal E_2,\ \|h\|=1}\Bigl(\sum_i\|x_i-w_ih\|^2\Bigr)^{1/2}.
\]
Thus $d_\cH(x)$ is the distance, in the product norm, to tuples with two
equal collinear vectors of length $1/\sqrt2$ (up to sign) and zeros
elsewhere. For real tuples $d_\cH=d$.

\begin{theorem}[Stability]\label{thm:stability}
There is an absolute constant $c_\cH>0$ such that, for every unit-energy
tuple in every Hilbert space,
\[
 c_\cH\,d_\cH(x)\ \le\ D(x)\ \le\ d_\cH(x).
\]
One may take $c_\cH=\min\{1/(200K),1/(4\sqrt2)\}$ with
$K=(2+24\sqrt2)/c_0+48$, where $c_0$ is any constant in the robust
Khintchine inequality $\E|S_a|\ge1/\sqrt2+c_0d(a)$. In particular equality
in \cref{thm:hilbert} for $f(t)=\min(t,1)$ holds only at two equal
collinear vectors.
\end{theorem}

\begin{proof}
\emph{Upper bound.} Couple $x$ with a nearest extremal tuple $y$ using the
same signs. Since $\min(\cdot,1)$ is $1$-Lipschitz and $D(y)=0$, Cauchy--Schwarz
gives $D(x)\le\E\|\sum_i\varepsilon_i(x_i-y_i)\|\le d_\cH(x)$.

\emph{Scalar lower bound.} For a real unit vector $a$, \cref{thm:stoploss}
gives $D(a)=\E|S_a|-\E(|S_a|-1)_+-\tfrac12\ge\E|S_a|-1/\sqrt2\ge c_0d(a)$, by
the robust Khintchine inequality of De, Diakonikolas and Servedio
\cite[Theorem~9]{DDS}.

\emph{Far from the extremal set.} Use the spectral decomposition of the
normalised Gram matrix $A=\sum_kp_kv_kv_k^{\mathsf T}$ from the proof of
\cref{thm:hilbert}, now with $s=1$. For fixed signs let
$y_k=\min(|v_k\cdot\varepsilon|,1)\in[0,1]$. Since
$\sqrt{\sum p_ky_k^2}-\sum p_ky_k=\Var_p(y)/(\sqrt{\sum p_ky_k^2}+\sum p_ky_k)
\ge\Var_p(y)/2$, inequality \eqref{eq:lift} and the scalar bound give
\begin{equation}\label{eq:variance}
 D(x)\ \ge\ c_0\sum_kp_kd(v_k)+\tfrac12\E\Var_p(y).
\end{equation}
Choose $w_k\in\mathcal E_2$ nearest to $v_k$, with $d(v_k)\le\sqrt2$, and
let $z_k=\min(|w_k\cdot\varepsilon|,1)$. Then $z_k$ is the indicator of a
parity event $\{\varepsilon_i\varepsilon_j=\sigma\}$ determined by the
rank-one projector $w_kw_k^{\mathsf T}$, and two different such events
disagree with probability at least $1/2$. Also
$\E(y_k-z_k)^2\le d(v_k)^2\le\sqrt2\,d(v_k)$. Group the indices $k$ by
their projectors, and let $p_*$ be the largest group mass, with projector
$P_*$. Drawing $k,l$ independently with weights $p$,
\[
 1-p_*\le2\E(z_k-z_l)^2\le6\bigl[\E(z_k-y_k)^2+\E(y_k-y_l)^2+\E(y_l-z_l)^2\bigr]
 \le\Bigl(\frac{12\sqrt2}{c_0}+24\Bigr)D(x).
\]
For unit vectors, $\|uu^{\mathsf T}-vv^{\mathsf T}\|_1\le2\|u-v\|$ in trace
norm, and distinct projectors are at most $2$ apart, so
\begin{equation}\label{eq:gram}
 \|A-P_*\|_1\le2\sum_kp_kd(v_k)+2(1-p_*)\le K\,D(x).
\end{equation}
Minimising over $h$ gives
$d_\cH(x)^2=2-2\max_{w\in\mathcal E_2}\sqrt{w^{\mathsf T}Aw}$, and for the
$w$ defining $P_*$ we have $w^{\mathsf T}Aw\ge1-\|A-P_*\|_1$. Hence
$d_\cH(x)^2\le2KD(x)$, which gives $D(x)\ge d_\cH(x)/(200K)$ whenever
$d_\cH(x)\ge1/100$.

\emph{Near the extremal set.} Suppose $d_\cH(x)<1/100$. Permute and change
signs so that the minimising pair is $x_1,x_2$ with equal positive weights.
Put $u=x_1+x_2$, $v=x_1-x_2$, $Z=\sum_{i\ge3}\varepsilon_ix_i$,
$b=\E\|Z\|^2$ and $L^2=\|v\|^2+b$. Then
\[
 d_\cH(x)^2=2-\sqrt2\|u\|,\qquad \|u\|^2=2-2b-\|v\|^2,\qquad
 d_\cH(x)/\sqrt2\le L\le\sqrt2\,d_\cH(x),
\]
the last by writing $d_\cH(x)^2=2(1-\sqrt{1-b-\|v\|^2/2})$. So $L<1/50$
and $\|u\|>7/5$. Conditioning on the first two signs and using
$\min(\|u+Z\|,1)\ge1-1_{\{\|u+Z\|<1\}}$,
\[
 2D(x)\ \ge\ \E\min(\|v+Z\|,1)-\PP(\|u+Z\|<1).
\]
By the spectral argument without clipping and Szarek's inequality, every
signed Hilbert sum has first moment at least $1/\sqrt2$ times the root of
its energy; with an independent sign in front of $v$ (allowed by symmetry
of $Z$) this gives $\E\|v+Z\|\ge L/\sqrt2$. The bound $(y-1)_+\le y^2/4$
costs at most $L^2/4$, and Chebyshev's inequality gives
$\PP(\|u+Z\|<1)\le b/(\|u\|-1)^2\le25L^2/4$. Hence
$D(x)\ge L/(2\sqrt2)-13L^2/4\ge L/4\ge d_\cH(x)/(4\sqrt2)$.
\end{proof}

\begin{remark}[Explicit constants]\label[remark]{rem:explicit}
Eskenazis, Nayar and Tkocz give the explicit robust constant
$c_0=8\cdot10^{-5}$ \cite[Theorem~8 and Section~5]{ENT}, which yields
$c_\cH>1.1\cdot10^{-8}$. Fang and Wang's preprint \cite{FangWang} gives
$c_0=1/20$, which would yield $c_\cH>6.5\cdot10^{-6}$. The scalar estimate
$c_0d(a)\le D(a)\le d(a)$ holds with the same $c_0$. These constants are far
from optimal; the point of \cref{thm:stability} is the linear order, which
is sharp because $D\le d_\cH$.
\end{remark}

%======================================================================
\section{Three equal weights}\label{sec:triadic}

This section extends Fang and Huang's inequality for three equal weights to
all weight vectors and proves the stability needed in \cref{sec:faces}.

Fang and Huang proved \cite[Lemma~2.4]{FangHuang} that if
$\sum_ia_i^2=3$ and $|a_i|\le1$ for all $i$, then
$\E(|S_a|-2)_+\le1/4$, with equality exactly for three weights of magnitude
one. Their proof uses the majorant \cite[Lemma~2.3]{FangHuang}
\begin{equation}\label{eq:FHmajorant}
 \Lambda(u)=u^2-1+\frac4\pi\cos\frac{\pi u}{2}-\frac{2}{3\pi}\cos^2\frac{\pi u}{2}
 \ \ge\ 8(|u|-2)_+,
\end{equation}
with equality exactly at $|u|\in\{1,3,5\}$, together with a cosine-product
inequality \cite[Lemma~2.2]{FangHuang}.

\begin{theorem}\label{thm:triadic}
There are absolute constants $c_3>0$ and $\eta>0$ such that, for every unit
vector $a$ and every $t$ with $|t-2/\sqrt3|\le\eta$,
\[
 \E(|S_a|-t)_+\le\frac{\sqrt3-t}{4}-c_3\dist(a,\mathcal E_3)^2 .
\]
Equality in the bound without the distance term holds exactly on
$\mathcal E_3$, and the quadratic order is best possible.
\end{theorem}

The constants $c_3$ and $\eta$ come from a compactness argument and are not
explicit. We work with weights normalised by $\sum_ia_i^2=3$ and raw
threshold $u=\sqrt3\,t$, so the claim is
$\E(|S_a|-u)_+\le(3-u)/4-c\,\delta(a)^2$ with $\delta(a)=\dist(a,\sqrt3\mathcal E_3)$.

\begin{proof}
\emph{Large weights.} Let $m$ be the largest magnitude and suppose $m>1$;
then $m\le\sqrt3<2$. Write $S_a=m\varepsilon+R$. Symmetry gives
$\E(|m\varepsilon+R|-2)_+=\tfrac12\E[(|R|-(2-m))_++(|R|-(2+m))_+]$. For
$x\ge0$ the bracket is at most $x^2/[4(2-m)]$: its two affine pieces are
$x-(2-m)$ and $2x-4$, with quadratic majorants $x^2/[4(2-m)]$ and $x^2/4$.
Since $\E R^2=3-m^2$,
\begin{equation}\label{eq:triple-large}
 \E(|S_a|-2)_+\le\frac{3-m^2}{8(2-m)}=\frac14-\frac{(m-1)^2}{8(2-m)}<\frac14 .
\end{equation}
For $m\le1$, Fang and Huang's lemma gives the inequality and the equality
case.

\emph{Compactness.} Consider a sequence of weight vectors (ordered
decreasingly in magnitude, padded with zeros) whose deficit at $u=2$ tends
to zero. By \eqref{eq:triple-large}, the largest magnitude has limit at most
one. Pass to coordinatewise limits $b_i$; then $\sum_ib_i^2\le3$ and the
sums converge in law to
$T=\sum_ib_i\varepsilon_i+\sigma G$ with $\sigma^2=3-\sum_ib_i^2$ and $G$
independent standard normal. (The $i$-th ordered weight satisfies
$a_i^2\le3/i$, so the fourth powers beyond index $m$ sum to at most
$9/(m+1)$ uniformly, and the tail is asymptotically Gaussian.) Fourth
moments are at most $27$, so $\E(|T|-2)_+=1/4$ in the limit. Approximating
$T$ by the first $m$ weights together with many equal small weights of the
right total energy (all of magnitude at most one) and applying
\eqref{eq:FHmajorant} and Fang and Huang's argument to the approximants
gives $\E\Lambda(T)\le2$. Since $\Lambda(T)-8(|T|-2)_+\ge0$ has expectation
at most zero, $T$ is supported on $\{\pm1,\pm3,\pm5\}$. A Gaussian component
would make $T$ continuous, so $\sigma=0$. Infinitely many nonzero $b_i$ would
give infinitely many support points, since the first $m$ positive weights
already produce $m+1$ distinct sums. So $T$ is a finite Rademacher sum with
weights at most one and zero deficit, and Fang and Huang's equality case
shows that it has exactly three weights of magnitude one. Since
$\sum_ib_i^2=3$, coordinatewise convergence upgrades to convergence in
$\ell^2$. Hence the deficit at $u=2$ is bounded below by a positive
constant outside every fixed neighbourhood of the equality set.

\emph{Near the equality set.} After ordering and sign changes put
$\delta^2=\sum_{i\le3}(a_i-1)^2+\sum_{i>3}a_i^2$, $s=a_1+a_2+a_3$ and
$v=\sum_{i>3}a_i^2$, so that $\delta^2=6-2s$ and $v\le\delta^2$. Suppose
$\delta\le1/32$ and $u\in[7/4,9/4]$. The eight signed sums $c$ of the first
three terms lie within $3/32$ of $\pm1$ or $\pm3$, hence at distance more
than $1/2$ from the kinks $\pm u$ of $\ell_u(x)=(|x|-u)_+$; two of them lie
outside $[-u,u]$ and six inside, so their mean excess is $(s-u)/4$. For the
independent remainder $R$, convexity and the distance to the kinks give
\[
 0\le\ell_u(c+R)-\ell_u(c)-\ell_u'(c)R\le2|R|1_{\{|R|\ge1/2\}}\le16R^4 .
\]
As $\E R=0$ and $\E R^4\le3v^2$, the deficit from $(3-u)/4$ is at least
$(3-s)/4-48v^2\ge\delta^2/8-48\delta^4\ge\delta^2/16$.

\emph{Conclusion.} Outside the neighbourhood $\delta\le1/32$ compactness
gives a positive gap at $u=2$. The deficit relative to $(3-u)/4$ is
$3/4$-Lipschitz in $u$, so a small fixed interval around $u=2$ keeps half
this gap; as $\delta$ is bounded, the gap is a multiple of $\delta^2$.
Rescaling by $1/\sqrt3$ gives the theorem. Perturbing only the first three
weights on the sphere gives deficit exactly $\delta^2/8$ at $u=2$, so the
quadratic order is sharp.
\end{proof}

%======================================================================
\section{Sets of large first-level weight near density one quarter}\label{sec:cube}

This section applies \cref{thm:stoploss} to the discrete cube. We first
prove a localisation estimate and the leading term of the local profile
without any recent input, and then use Yu's small-density theorem to obtain
the finer structure.

\paragraph{Set-up.}
For $A\subseteq\{-1,1\}^n$ with density $\mu(A)$ and centroid vector
$b_A=\E[\varepsilon1_A]$, the number $\|b_A\|^2$ is the first-level Fourier
weight of $1_A$. If $X,Y$ are independent and uniform on $A$, their mean
Hamming distance is $n/2-\|b_A\|^2/(2\mu(A)^2)$ \cite[Lemma~2.1]{FangHuang},
\cite[equation~(4)]{YuLevel}, so maximising $\|b_A\|^2$ at fixed density is
the Ahlswede--Katona average-distance problem \cite{AhlswedeKatona}. We use
the dimension-free profile
\[
 W(\alpha)=\lim_{\eta\downarrow0}\ \sup\bigl\{\|b_A\|^2:\ n\ge1,\
 A\subseteq\{-1,1\}^n,\ |\mu(A)-\alpha|<\eta\bigr\}.
\]
At dyadic densities $W$ equals the supremum at exact density, and $W$ is the
continuous extension used by Yu \cite[Section~II-D]{YuLevel}. The values
$W(1/4)=1/8$ \cite{FuWeiYeung} and $W(1/8)=3/64$ \cite{FangHuang} are
known, with $2$-faces and $3$-faces as the extremal sets. All logarithms are
natural, and $h\downarrow0$ is taken after the dimension-free supremum.

We use the following signed form of the first-level inequality (Chang's
lemma; compare \cite[Theorem~1]{YuLevel}): for an event $B$ of mass $p$
and a sign-valued function $\xi$ on $B$, the vector
$v=\E[\xi\varepsilon1_B]$ satisfies
\begin{equation}\label{eq:signedchang}
 \|v\|\le p\sqrt{2\log(2/p)} .
\end{equation}
For a unit vector $u$, $u\cdot v\le p\,\E[|S_u|\mid B]$, and Jensen's
inequality with $\E e^{\lambda|S_u|}\le2e^{\lambda^2/2}$ gives
$\E[|S_u|\mid B]\le(\log(2/p)+\lambda^2/2)/\lambda$; take
$\lambda=\sqrt{2\log(2/p)}$.

\begin{proposition}[Localisation]\label[proposition]{prop:cubestable}
Let $M=(\sqrt2-1)/4$ and $\Delta(A)=M-\|b_A\|+\mu(A)\ge0$. For $n\ge2$,
\[
 c\min_Q\mu(A\triangle Q)\ \le\ \Delta(A)\ \le\ C\min_Q\mu(A\triangle Q),
\]
where $Q$ ranges over $2$-faces, $C=1+\sqrt2$ and
$1/c=2/\gamma+8000/\gamma^2$ with $\gamma=\sqrt2-1$.
\end{proposition}

\begin{proof}
For a unit vector $a$, symmetry and \cref{thm:stoploss} give
$\E[(S_a-1)1_A]\le\E(S_a-1)_+=\tfrac12\E(|S_a|-1)_+\le M$. Choose $a$ in the
direction of $b_A$ (arbitrarily if $b_A=0$), so that $\E[S_a1_A]=\|b_A\|$;
this shows $\Delta(A)\ge0$, and $\Delta(A)$ is the sum of the two
nonnegative slacks $\sigma_1=\E(S_a-1)_+-\E[(S_a-1)1_A]$ and
$\sigma_2=M-\E(S_a-1)_+$. By \eqref{eq:stab1}, $\sigma_2\ge d(a)^2/2000$.
Let $w\in\mathcal E_2$ be nearest to $a$; then $Q=\{S_w>1\}$ is a $2$-face,
and $|S_w-1|\ge\gamma$ everywhere. On the event $|S_a-S_w|<\gamma/2$, each
point of $A\triangle Q$ contributes at least $\gamma/2$ to $\sigma_1$
(through $(S_a-1)_+$ on $Q\setminus A$ and $(S_a-1)_-$ on $A\setminus Q$).
Chebyshev's inequality bounds the remaining event by
$4\|a-w\|^2/\gamma^2$. Hence
$\mu(A\triangle Q)\le2\sigma_1/\gamma+8000\sigma_2/\gamma^2\le\Delta(A)/c$.
Conversely, testing the direction $w$ of any $2$-face $Q$ gives
$\Delta(A)\le\E[(S_w-1)(1_Q-1_A)]\le(1+\sqrt2)\mu(A\triangle Q)$, since
$\E[(S_w-1)1_Q]=M$ and $|S_w-1|\le1+\sqrt2$.
\end{proof}

\begin{theorem}[Leading term]\label{thm:cusp}
As $h\downarrow0$,
\[
 W\bigl(\tfrac14\pm h\bigr)=2\min\bigl(\tfrac14,\tfrac14\pm h\bigr)^2+(2+o(1))\,h^2\log(1/h).
\]
In particular $W$ has left derivative $1$ and right derivative $0$ at $1/4$.
This result does not use \cite{YuLevel}.
\end{theorem}

\begin{proof}
\emph{Lower bound.} Let $Q=\{\varepsilon_1=\varepsilon_2=1\}$ and
$w=(e_1+e_2)/\sqrt2$. Below $1/4$ remove mass $h$ from $Q$; above $1/4$ add
mass $h$ inside the neighbouring face $\{\varepsilon_1=1,\varepsilon_2=-1\}$.
In both cases choose the removed or added points by thresholding the
normalised sum $T_m$ of $m$ further coordinates, so that the conditional mass
in the face is close to $4h$. The projection of $b_A$ on $w$ is
$p_0=\sqrt2\min(\alpha,1/4)$ with $\alpha=1/4\pm h$. By the central limit
theorem and uniform integrability (second moments equal one and the normal
law has no atoms), the component along the common direction of the $m$
coordinates tends to $I(4h)/4$ as $m\to\infty$. Hence
$W(\alpha)\ge p_0^2+I(4h)^2/16=p_0^2+(2+o(1))h^2\log(1/h)$, by
$I(p)\sim p\sqrt{2\log(1/p)}$.

\emph{Upper bound.} Let $A$ have density $1/4\pm h$, and assume
$\|b_A\|\ge p_0$ (otherwise there is nothing to prove). Then $\Delta(A)\le h$
above $1/4$ and $\Delta(A)\le(\sqrt2-1)h$ below it, so by
\cref{prop:cubestable} and symmetry there is a $2$-face, say $Q$ as above,
with $\delta=\mu(A\triangle Q)=O(h)$ uniformly in $n$. Put
$r=\mu(Q\setminus A)$, $s=\mu(A\setminus Q)$ and $z=\min(r,s)$, so that
$\delta=h+2z$. The projection $p$ of $b_A$ on $w$ satisfies
$p\le p_0-\sqrt2z$: removal from $Q$ costs $\sqrt2$ per unit of mass and
additions outside $Q$ have nonpositive projection. Also $p=p_0+O(h)$. The
component of $b_A$ orthogonal to $w$ is that of $b_A-b_Q$, of norm at most
$\delta\sqrt{2\log(2/\delta)}$ by \eqref{eq:signedchang}. From
$\|b_A\|^2\ge p_0^2$ we get $(p_0-p)(p_0+p)\le2\delta^2\log(2/\delta)$,
hence $z=O(h^2\log(1/h))$ and $\delta=h+O(h^2\log(1/h))$. Therefore
$\|b_A\|^2\le p_0^2+2\delta^2\log(2/\delta)=p_0^2+(2+o(1))h^2\log(1/h)$.
The bounds are uniform in $n$ and pass to $W$.
\end{proof}

\subsection{Finer structure, using Yu's small-density theorem}

The rest of this section, and the exact formulae of \cref{sec:faces}, use
the following theorem of Yu. Yu proves that there is a universal $a_{\mathrm Y}>0$
such that every Boolean function $g$ on any finite cube with mean
$p\in(0,a_{\mathrm Y}]$ satisfies $\sum_i\hat g_i^2<I(p)^2$, and that
$W(p)=I(p)^2$ for $0<p\le a_{\mathrm Y}$ \cite[Theorem~5, version~2]{YuLevel}. In our
notation: every set $B$ of density $p\in(0,a_{\mathrm Y}]$ has $\|b_B\|<I(p)$, and
by complementation every set of density $1-p$ has $\|b_B\|<I(p)$. The
threshold $a_{\mathrm Y}$ is not explicit, so the neighbourhoods below have no
specified width.

\begin{corollary}[Exact formula; uses \cite{YuLevel}]\label[corollary]{cor:fourfibres}
For $|\alpha-1/4|$ sufficiently small,
\[
 W(\alpha)=\frac1{16}\max\Bigl\{\Bigl\|\sum_{e\in\{-1,1\}^2}e\,p_e\Bigr\|^2
 +\Bigl(\sum_{e\in\{-1,1\}^2}I(p_e)\Bigr)^2:\ 0\le p_e\le1,\ \sum_ep_e=4\alpha\Bigr\}.
\]
The maximum is over three free real parameters.
\end{corollary}

\begin{proof}
\emph{Lower bound.} Condition on the first two signs $e$. In the fibre over
$e$, take an upper threshold set of the \emph{same} normalised sum of the
remaining coordinates with conditional mass close to $p_e$. As in
\cref{thm:cusp}, the first two centroid coordinates tend to
$\sum_eep_e/4$ and the component along the common direction tends to
$\sum_eI(p_e)/4$.

\emph{Upper bound.} Take sets with densities tending to $\alpha$ and
weights tending to $W(\alpha)$. By the lower bound, their weights are at
least $2\min(\alpha,1/4)^2-o(1)$, so \cref{prop:cubestable} places them
within $O(|\alpha-1/4|)+o(1)$ of $2$-faces. Use the two coordinates of the
face. One conditional density is close to one and the other three are
close to zero; each $\min(p_e,1-p_e)$ is at most four times the symmetric
difference. Choose the neighbourhood so that these lie in $(0,a_{\mathrm Y}]$. By
Yu's theorem each conditional centroid (in the remaining coordinates) has
norm at most $I(p_e)$, and the triangle inequality bounds the remaining
part of $b_A$ by $\sum_eI(p_e)/4$. Compactness and continuity of $I$ pass to
a limiting tuple.
\end{proof}

The tuples realising the maximum can be taken \emph{localised}, that is,
within $O(|\alpha-1/4|)$ of $(1,0,0,0)$ in the order $++,+-,-+,--$ up to
symmetry. Indeed, in the lower-bound construction every coordinate of
$b_A$ beyond the first two is $O(m^{-1/2})$, so for large $m$ the nearest
direction in $\mathcal E_2$ uses the first two coordinates, and
\cref{prop:cubestable} then applies to the face they define.

Write $H_2(\lambda)=-\lambda\log\lambda-(1-\lambda)\log(1-\lambda)$ for the
binary entropy and put
\[
 \kappa_2=\max_{0\le\lambda\le1}\bigl\{2H_2(\lambda)+2(2\lambda-1)^2\bigr\}
 =2.0393421359\ldots .
\]

\begin{theorem}[Next-order terms; uses \cite{YuLevel}]\label{thm:refined}
As $h\downarrow0$,
\begin{align*}
 W(\tfrac14-h)&=2(\tfrac14-h)^2+\tfrac1{16}I(4h)^2+o(h^2),\\
 W(\tfrac14+h)&=\tfrac18+\tfrac1{16}I(4h)^2+\kappa_2h^2+o(h^2),
\end{align*}
where
$\tfrac1{16}I(4h)^2=h^2\{2\log(1/h)-\log\log(1/h)+2-\log(64\pi)+o(1)\}$.
Above $1/4$, maximising tuples put essentially all added mass in the two
faces adjacent to the main face, in proportions converging to
$\lambda_*=(1+m_*)/2=0.9787520120\ldots$ and $1-\lambda_*$, where $m_*$ is
the unique root of $m=\tanh(2m)$ in $(0,1)$, up to exchanging the two faces.
\end{theorem}

The two expansions share the $h^2\log(1/h)$ and $h^2\log\log(1/h)$ terms
and differ in the constant multiple of $h^2$; above $1/4$ the gain
$\kappa_2h^2$ comes from splitting the added mass unequally between two
faces. The proportions describe maximisers of the reduced problem in
\cref{cor:fourfibres}; whether every near-optimal \emph{set} has these
proportions is left open (\cref{sec:open}).

\begin{proof}
By \cref{cor:fourfibres} and the remark after it, write $q=4h$,
$L=\log(1/q)$ and take localised conditional masses
$(1-r,x,y,z)$ (order $++,+-,-+,--$) with $r,x,y,z=O(q)$. The objective is
\begin{equation}\label{eq:reduced-objective}
 \frac{(1-r-z)^2+(x-y)^2}{8}+\frac{[I(r)+I(x)+I(y)+I(z)]^2}{16}.
\end{equation}

\emph{Removing the wasteful allocations.} Above $1/4$ set $w=r+z$,
$s=x+y=q+r-z$ and $\lambda=x/s$ (any $\lambda$ if $s=0$), and compare with
$(r,x,y,z)=(0,q\lambda,q(1-\lambda),0)$. The first term of
\eqref{eq:reduced-objective} is larger for the comparison tuple by at least
$w/4-O(qw)$. The adjacent Gaussian terms of the original tuple exceed those
of the comparison tuple by at most $Cw\sqrt L$ when $s>q$ (use
$I'(p)=-\Phi^{-1}(p)$ and
$\lambda\sqrt{\log(1/(q\lambda))}+(1-\lambda)\sqrt{\log(1/(q(1-\lambda)))}\le\sqrt{L+H_2(\lambda)}$),
and are smaller when $s\le q$. The terms $I(r)+I(z)$ are at most
$2I(w/2)\le Cw\sqrt{\log(2/w)}$ by concavity of $I$. All Gaussian sums are
$O(q\sqrt L)$ and $w=O(q)$. Expanding the square, the objective of the
original tuple exceeds that of the comparison tuple by at most
\begin{equation}\label{eq:wrong-fibres}
 -\frac w8+Cq\sqrt L\,w\sqrt{\log(2/w)}
\end{equation}
for small $q$, after absorbing terms of order $qwL$ (including the square
of $I(r)+I(z)$, since $w\log(2/w)=O(qL)$ for $w=O(q)$). If $w\ge q^3$, the
second term is $O(qLw)$, so \eqref{eq:wrong-fibres} is negative and the
comparison tuple is strictly better. If $w<q^3$, \eqref{eq:wrong-fibres}
is $O(q^4L)=o(q^2)$. Therefore
\begin{equation}\label{eq:adjacent-reduction}
 W(\tfrac14+h)=\frac18+\max_{0\le\lambda\le1}\Bigl\{\frac{q^2(2\lambda-1)^2}{8}
 +\frac{[I(q\lambda)+I(q(1-\lambda))]^2}{16}\Bigr\}+o(q^2),
\end{equation}
with all constants independent of $\lambda$. Below $1/4$, $r=q+x+y+z$;
put $w=x+y+z$ and compare with $(q,0,0,0)$. The first term again loses at
least $w/4-O(qw)$, $I(q+w)-I(q)\le Cw\sqrt L$ and
$I(x)+I(y)+I(z)\le3I(w/3)$, so the same estimate proves the first
expansion.

\emph{Entropy expansion.} We claim that, uniformly in $\lambda\in[0,1]$,
\begin{equation}\label{eq:uniform-entropy}
 [I(q\lambda)+I(q(1-\lambda))]^2=I(q)^2+2q^2H_2(\lambda)+o(q^2).
\end{equation}
Put $R(p)=I(p)/p$ and $\zeta_p=-\Phi^{-1}(p)$. Integration by parts in the
Gaussian tail gives $0<R(p)-\zeta_p<1/\zeta_p$ and
$R(p)-\zeta_p\sim1/\zeta_p$; also $dR/d\log(1/p)=R(p)-\zeta_p$ and
$\zeta_q\sim\sqrt{2L}$. For fixed $\lambda>0$ it follows that
$\sqrt{2L}\,\lambda\{R(q\lambda)-R(q)\}\to\lambda\log(1/\lambda)$, and for
all $\lambda\in(0,1]$ this expression lies between $0$ and
$(\sqrt{2L}/\zeta_q)\lambda\log(1/\lambda)$. So the convergence is uniform
on $[0,1]$. Adding the two expressions, squaring and using
$I(q)/q\sim\sqrt{2L}$ gives \eqref{eq:uniform-entropy}. Substituting in
\eqref{eq:adjacent-reduction} with $q=4h$ gives the second expansion.
Two further integrations by parts give
$R(p)^2=2\log(1/p)-\log\log(1/p)-\log(4\pi)+2+o(1)$, which yields the
displayed form of $I(4h)^2/16$.

\emph{The allocation.} With $m=2\lambda-1$, the derivative of
$2H_2((1+m)/2)+2m^2$ is $4m-2\operatorname{artanh}m$. Since
$\operatorname{artanh}(m)/m$ increases from $1$ to $\infty$ on $(0,1)$,
there is exactly one positive stationary point, given by
$\operatorname{artanh}m=2m$, equivalently $m=\tanh(2m)$; the centre is a
strict local minimum and the derivative is negative near $m=1$. So the two
symmetric points are the global maxima, and uniformity in
\eqref{eq:uniform-entropy} transfers this to maximising tuples.
\end{proof}

%======================================================================
\section{Exposed subcubes and density one eighth}\label{sec:faces}

This section shows that the argument of \cref{sec:cube} works near
density $2^{-k}$ whenever $k$ equal weights satisfy a stable stop-loss
inequality, and applies it with $k=3$ via \cref{thm:triadic}.

\begin{theorem}[Transfer; uses \cite{YuLevel}]\label{thm:faces}
Let $k\ge2$ and $a_0=2^{-k}$. Suppose that some
$t\in((k-2)/\sqrt k,\sqrt k)$ and $c>0$ satisfy, for all unit vectors $a$,
\begin{equation}\label{eq:exposed}
 a_0(\sqrt k-t)-\E(S_a-t)_+\ \ge\ c\,\dist(a,\mathcal E_k)^2 .
\end{equation}
Then:
\begin{enumerate}[label=(\roman*)]
\item For all $\alpha$ sufficiently close to $a_0$,
\begin{equation}\label{eq:k-fibres}
 W(\alpha)=4^{-k}\max\Bigl\{\Bigl\|\sum_{e\in\{-1,1\}^k}e\,p_e\Bigr\|^2
 +\Bigl(\sum_eI(p_e)\Bigr)^2:\ 0\le p_e\le1,\ \sum_ep_e=2^k\alpha\Bigr\},
\end{equation}
and for $0<|\alpha-a_0|$ sufficiently small no set in finite dimension of
exact density $\alpha$ has $\|b_A\|^2=W(\alpha)$.
\item With $q=2^kh$, as $h\downarrow0$,
\begin{align}
 W(a_0-h)&=k(a_0-h)^2+4^{-k}I(q)^2+o(h^2),\notag\\
 W(a_0+h)&=ka_0^2+2a_0(k-2)h+4^{-k}I(q)^2+\kappa_kh^2+o(h^2),\label{eq:k-constant}
\end{align}
where
\[
 \kappa_k=k-4+\max_{\lambda\in\mathcal P_k}\Bigl\{4\sum_{i=1}^k\lambda_i^2+2\Ent(\lambda)\Bigr\}
\]
and $4^{-k}I(2^kh)^2=h^2\{2\log(1/h)-\log\log(1/h)+2-\log(4^{k+1}\pi)+o(1)\}$.
\end{enumerate}
\end{theorem}

\begin{proof}
Constants $C_k$ and implied constants may depend on $k$, $t$ and $c$, but
not on the dimension.

\emph{Localisation.} Put $M=a_0(\sqrt k-t)$ and
$\Delta=M-\|b_A\|+t\mu(A)$. Testing the direction of $b_A$ splits $\Delta$
into the stop-loss deficit and a classification loss, as in
\cref{prop:cubestable}. If $w\in\mathcal E_k$ is nearest to this direction,
$Q=\{S_w>t\}$ is the $k$-face determined by $w$, because $S_w$ takes the
values $(k-2j)/\sqrt k$ and $t$ lies strictly between the two largest. Every
value of $S_w$ is at least
$\gamma=\min(\sqrt k-t,\ t-(k-2)/\sqrt k)>0$ from $t$, and the argument of
\cref{prop:cubestable} with \eqref{eq:exposed} gives
$\mu(A\triangle Q)\le C_k\Delta$; testing $w$ gives the converse bound.

\emph{Exact formula.} Removing mass $h$ from a $k$-face, or adding it to
the $k$ neighbouring faces, gives centroid norms at least
$\sqrt k(a_0-h)$ below $a_0$ and $\sqrt ka_0+(k-2)h/\sqrt k$ above it. So a
maximising sequence at nearby density has $\Delta=O(h)+o(1)$ and lies within
$O(h)+o(1)$ of a $k$-face. Decompose by the $k$ fixed signs into $2^k$
fibres: one conditional density is near one and the others near zero. For
a small enough neighbourhood all minority densities lie in Yu's range, and
the argument of \cref{cor:fourfibres} gives \eqref{eq:k-fibres}, with
localised maximising tuples as explained after it.

\emph{Nonattainment.} Suppose a finite-dimensional set $A$ of exact density
$\alpha\ne a_0$ had $\|b_A\|^2=W(\alpha)$. Its weight is at least that of
the construction above, so $\Delta(A)=O(h)$ and $A$ is localised near a
$k$-face. As $2^k\alpha$ is not an integer near $1$, some fibre has
conditional density strictly between $0$ and $1$, and lying in $(0,a_{\mathrm Y}]$
after complementation. Yu's strict finite-dimensional inequality makes that
fibre's centroid strictly shorter than $I(p_e)$, so $\|b_A\|^2$ is strictly
below the objective of its own tuple, which is at most $W(\alpha)$.

\emph{Next order.} Index the faces of $\{-1,1\}^k$ by the set $J$ of
negative coordinates. A localised tuple has $p_\varnothing=1-r$ and
$p_J=x_J$ for $J\ne\varnothing$, with $r,x_J=O(q)$, and the objective is
\[
 4^{-k}\Bigl\{\Bigl\|(1-r)\mathbf1+\sum_{J\ne\varnothing}x_Je_J\Bigr\|^2
 +\Bigl[I(r)+\sum_{J\ne\varnothing}I(x_J)\Bigr]^2\Bigr\},
\]
where $e_J$ is the sign vector of face $J$. Above $a_0$ the constraint is
$\sum_Jx_J-r=q$. Compare with the tuple that puts mass $q\lambda_i$ on the
singleton faces $J=\{i\}$, in the proportions of the current singleton
masses, and nothing elsewhere. With $w=r+\sum_{|J|\ge2}x_J$, the first
squared norm of the comparison tuple exceeds that of the original by at
least $4w-O(qw)$: the inner product of $e_J$ with $\mathbf1$ is $k-2|J|$,
so non-singleton faces and removal lose at least $4$ per unit of mass
relative to singleton faces. The Gaussian terms change by at most
$C_kq\sqrt L\,w\sqrt{\log(C_k/w)}$, using
$\sum_i\lambda_i\sqrt{\log(1/(q\lambda_i))}\le\sqrt{L+\Ent(\lambda)}$ and
$\Ent(\lambda)\le\log k$. The splitting at $w=q^3$ from \cref{thm:refined}
shows that the original tuple gains at most $o(q^2)$. Below $a_0$ compare
with pure removal $r=q$; with $w=\sum_{J\ne\varnothing}x_J$ the first term
loses at least $4\sum_J|J|x_J-O(qw)\ge4w-O(qw)$ and the same bound applies.
For the singleton tuple the retained centroid is
$a_0\mathbf1+h(\mathbf1-2\lambda)$, with squared norm
$ka_0^2+2a_0(k-2)h+h^2(k-4+4\sum_i\lambda_i^2)$, and the uniform
expansion $(\sum_iI(q\lambda_i))^2=I(q)^2+2q^2\Ent(\lambda)+o(q^2)$
follows from the proof of \eqref{eq:uniform-entropy} summed over the $k$
terms. This gives \eqref{eq:k-constant}; pure removal gives the lower
expansion, and the Gaussian asymptotic is as before.
\end{proof}

Hypothesis \eqref{eq:exposed} holds for $k=2$ (with $t=1$) by
\cref{thm:stoploss}, since $\E(S_a-t)_+=\tfrac12\E(|S_a|-t)_+$, and for
$k=3$ (with $t=2/\sqrt3$) by \cref{thm:triadic}. For $k=2$,
\cref{thm:faces} recovers \cref{cor:fourfibres,thm:refined}.

\begin{corollary}[Density one eighth; uses \cite{YuLevel}]\label[corollary]{cor:eighth}
As $h\downarrow0$,
\begin{align*}
 W(\tfrac18-h)&=3(\tfrac18-h)^2+\tfrac1{64}I(8h)^2+o(h^2),\\
 W(\tfrac18+h)&=\tfrac3{64}+\tfrac h4+\tfrac1{64}I(8h)^2+\kappa_3h^2+o(h^2),
\end{align*}
with $\kappa_3=3.0815948865\ldots$. The left and right derivatives of $W$
at $1/8$ are $3/4$ and $1/4$. Above $1/8$, maximising tuples split the added
mass among the three neighbouring faces in proportions converging to a
permutation of $(p_*,(1-p_*)/2,(1-p_*)/2)$, where $p_*=0.9539302689\ldots$
is the unique root in $(1/3,1)$ of $\log(2p/(1-p))=6p-2$.
\end{corollary}

\begin{proof}
Apply \cref{thm:faces} with $k=3$. For the maximisation in $\kappa_3$:
entropy rules out a zero coordinate at a maximum; the stationarity
equations make $4\lambda_i-\log\lambda_i$ constant, so there are at most
two distinct values; two coordinates above $1/4$ give a positive second
variation when mass is exchanged between them; and the uniform vector is a
strict local minimum. So a maximiser has the form
$(p,(1-p)/2,(1-p)/2)$. Along this line the derivative is
$-2[\log(2p/(1-p))-6p+2]$; the bracket vanishes at $p=1/3$, decreases and
then increases to $+\infty$ on $(1/3,1)$, so it has exactly one zero there,
which is a strict maximum. Substitution gives $\kappa_3$.
\end{proof}

\begin{remark}[Scope of the transfer theorem]\label[remark]{rem:scope}
Hypothesis \eqref{eq:exposed} requires $k$ equal weights to be the unique
maximisers of the stop-loss transform at some threshold. Among
equal-weight vectors, numerical evaluation (\cref{sec:computations}) shows
that two equal weights win for $t$ between $2-\sqrt2$ and $2\sqrt2-\sqrt3$,
three for $t$ up to $2\sqrt3-2\approx1.464$, four up to about $1.72$, and that for
$k=5$ and $k=6$ another equal-weight vector or the Gaussian always does
strictly better on the whole admissible interval. On an evaluation grid of
step below $5\cdot10^{-4}$ the margin is at least $10^{-3}$; since each
stop-loss transform is $1$-Lipschitz in $t$, the difference of two of them is
$2$-Lipschitz, so the margin stays positive between grid points (up to the
accuracy of the 30-digit evaluation). So the hypothesis fails
for $k=5,6$, and \cref{thm:faces} is expected to apply only for
$k\in\{2,3,4\}$. The case $k=4$ is Fang and Huang's
Conjecture~3.1 \cite{FangHuang} (with stability), which would settle
density $1/16$.
\end{remark}

\begin{remark}[The Potts connection]\label[remark]{rem:potts}
The maximisation defining $\kappa_k$ is, up to the factor $-2$ and the
additive constant $k-4$, the zero-field Curie--Weiss--Potts free-energy
minimisation at inverse temperature $\beta=4$: in the convention of
Gandolfo, Ruiz and Wouts \cite[equation~(2.3)]{Potts} the free energy is
$\sum_i\lambda_i\log\lambda_i-\tfrac\beta2\sum_i\lambda_i^2$. At this
temperature, for $k=2$ and $k=3$, the minimisers have one large coordinate
(\cref{thm:refined,cor:eighth}), which is why the added mass prefers one
neighbouring face. We use only the elementary
analysis given above for $k=2,3$.
\end{remark}

%======================================================================
\section{Computations}\label{sec:computations}

This section describes the accompanying scripts and what each one checks.
The proofs do not depend on them.

\begin{itemize}[leftmargin=1.5em]
\item \texttt{verify.py} (standard library) checks \cref{thm:stoploss} at
 $t\in\{1/\sqrt2,3/4,1,201/200\}$ and \cref{thm:triadic}'s bound at
 $t=2/\sqrt3$, with exact equality cases, for all 1{,}000 nondecreasing
 integer weight tuples with entries $1,\ldots,4$ and length at most $10$,
 using integer arithmetic and signs in $\mathbb Q(\sqrt2)$. It also encloses
 $\E\min(\|S\|,1)$ for 923 tuples of integer vectors in $\R^2$.
\item \texttt{cube\_census.py} enumerates all $65{,}536$ subsets of
 $\{-1,1\}^4$: the maximal weight is $32/256$ at density $1/4$, attained
 exactly by the $24$ two-faces, and $12/256$ at density $1/8$, attained
 exactly by the $32$ three-faces (one-dimensional subcubes). The supporting
 deficits in this census are evaluated in floating point with tolerance
 $10^{-14}$. It recomputes $m_*$, $\lambda_*$,
 $\kappa_2$, $p_*$, $\kappa_3$ and the Gaussian constants to 80 digits.
\item \texttt{verify\_asymptotic.py} (SciPy) evaluates the reduced
 one-parameter problems of \cref{thm:refined,cor:eighth} for
 $q=10^{-3},\ldots,10^{-200}$; the corrections approach $\kappa_2$ and
 $\kappa_3$ slowly, with error of order $1/\log(1/q)$. This illustrates
 the limit; it does not test the reductions.
\item \texttt{check\_extensions.py} (mpmath) evaluates the margins in the
 proof of \cref{thm:stoploss}, the endpoints in \cref{prop:interval}, the
 constants of \cref{cor:tv,rem:explicit,prop:cubestable}, the evidence for
 \cref{conj:interval} and the equal-weight envelope in \cref{rem:scope}.
 For \cref{conj:interval} it checks all nondecreasing integer tuples with
 entries $1,\ldots,5$ and length at most $9$ (2{,}001 tuples) at 23
 thresholds inside $(2-\sqrt2,2\sqrt2-\sqrt3)$, and 3{,}000 seeded random
 unit vectors of length at most $12$; no vector exceeds the pair.
\end{itemize}
Numerical results that are not exact (the envelope, the random vectors and
the asymptotic optimiser) carry a declared tolerance and are evidence, not
proof.

%======================================================================
\section{Open problems}\label{sec:open}

This section lists the questions that the paper leaves open.

\begin{conjecture}\label[conjecture]{conj:interval}
For every $t\in[2-\sqrt2,\,2\sqrt2-\sqrt3]$ and every unit vector $a$,
$\E(|S_a|-t)_+\le(\sqrt2-t)/2$. By \cref{prop:interval} this would be the
largest such interval. By convexity in $t$ it suffices to prove the two
endpoint cases, at each of which two different weight vectors tie.
\end{conjecture}

\begin{enumerate}[leftmargin=2em]
\item \emph{Density $1/16$.} Prove Fang and Huang's Conjecture~3.1
 \cite{FangHuang}, ideally with the stability required in \eqref{eq:exposed};
 by \cref{thm:faces} this would give the local profile near $1/16$.
 A three-harmonic analogue of \eqref{eq:FHmajorant} cannot work: matching
 the contact conditions at $0,2,4$ and the Gaussian limit leads to
 incompatible constraints on the coefficients.
\item \emph{Independence from Yu's theorem.} Find a proof of
 \cref{cor:fourfibres} or \cref{thm:refined} that avoids
 \cite[Theorem~5]{YuLevel}, for example a direct stop-loss bound for the
 small fibres at densities of order $h$.
\item \emph{Effective constants.} Make the error term in \cref{thm:cusp}
 effective, and make the neighbourhoods in
 \cref{cor:fourfibres,thm:faces} and the constants $c_3,\eta$ in
 \cref{thm:triadic} explicit.
\item \emph{Allocation for near-optimal sets.} Show that every set near
 density $1/4+h$ whose weight is within $o(h^2)$ of $W$ splits its added
 mass between two neighbouring faces in proportions near
 $\lambda_*:1-\lambda_*$.
\end{enumerate}

\section*{Statement on the use of AI}
This paper was prepared with substantial assistance from AI systems, which
were used to develop and check proofs, run the computations and draft the
text. The results were checked by independent AI-assisted proof audits
and by the computations described in \cref{sec:computations}. The paper has
not been peer reviewed.

\begin{thebibliography}{99}\small
\bibitem{AhlswedeKatona}
R. Ahlswede and G. O. H. Katona (1977). Contributions to the geometry of
Hamming spaces. \emph{Discrete Mathematics}, 17(1), 1--22.
\href{https://doi.org/10.1016/0012-365X(77)90017-6}{doi:10.1016/0012-365X(77)90017-6}.
\bibitem{AlthoferSillke}
I. Alth\"ofer and T. Sillke (1992). An ``average distance'' inequality for
large subsets of the cube. \emph{Journal of Combinatorial Theory, Series B},
56(2), 296--301.
\href{https://doi.org/10.1016/0095-8956(92)90024-R}{doi:10.1016/0095-8956(92)90024-R}.
\bibitem{DDS}
A. De, I. Diakonikolas and R. A. Servedio (2016). A robust Khintchine
inequality, and algorithms for computing optimal constants in Fourier
analysis and high-dimensional geometry. \emph{SIAM Journal on Discrete
Mathematics}, 30(2), 1058--1094.
\href{https://doi.org/10.1137/130919143}{doi:10.1137/130919143}.
\bibitem{ENT}
A. Eskenazis, P. Nayar and T. Tkocz (2024). Resilience of cube slicing in
$\ell_p$. \emph{Duke Mathematical Journal}, 173(17), 3377--3412.
\href{https://doi.org/10.1215/00127094-2024-0004}{doi:10.1215/00127094-2024-0004}.
\bibitem{FangHuang}
Y. Fang and H. Huang (2026). Average pairwise Hamming distance of a subset
in a hypercube. \emph{arXiv}:2610.40486v1 [math.CO], 29 October 2026.
\url{https://arxiv.org/abs/2610.40486}.
\bibitem{FangWang}
X. Fang and T. Wang (2026). A two-regime Khintchine inequality and an
improved bound on the degree-1 Fourier weight for linear threshold
functions. \emph{arXiv}:2608.27908v2 [math.PR], 1 November 2026.
\url{https://arxiv.org/abs/2608.27908}.
\bibitem{FuWeiYeung}
F.-W. Fu, V. K. Wei and R. W. Yeung (2001). On the minimum average distance
of binary codes: linear programming approach. \emph{Discrete Applied
Mathematics}, 111(3), 263--281.
\href{https://doi.org/10.1016/S0166-218X(00)00284-5}{doi:10.1016/S0166-218X(00)00284-5}.
\bibitem{Potts}
D. Gandolfo, J. Ruiz and M. Wouts (2010). Limit theorems and coexistence
probabilities for the Curie--Weiss Potts model with an external field.
\emph{Stochastic Processes and their Applications}, 120(1), 84--104.
\href{https://doi.org/10.1016/j.spa.2009.10.011}{doi:10.1016/j.spa.2009.10.011}.
\bibitem{Haagerup}
U. Haagerup (1981). The best constants in the Khintchine inequality.
\emph{Studia Mathematica}, 70(3), 231--283.
\href{https://doi.org/10.4064/sm-70-3-231-283}{doi:10.4064/sm-70-3-231-283}.
\bibitem{HP}
L. Hollom and J. Portier (2025). Tight anti-concentration of Rademacher
sums. \emph{Random Structures \& Algorithms}, 67(1).
\href{https://doi.org/10.1002/rsa.70024}{doi:10.1002/rsa.70024}.
Preprint arXiv:2306.07811 (theorem numbers refer to this version).
\bibitem{Kontorovich}
A. Kontorovich (2024). Expectation comparison inequality for concave
function of symmetric random variables. \emph{MathOverflow}, question
478477, 9 September 2024. \url{https://mathoverflow.net/q/478477}.
\bibitem{KontorovichTV}
A. Kontorovich (2024). On the tensorization of the variational distance.
\emph{arXiv}:2409.10368v2 [math.PR], 2 October 2024.
\url{https://arxiv.org/abs/2409.10368}.
\bibitem{Lowther}
G. Lowther (2024). Answer to ``Expectation comparison inequality for concave
function of symmetric random variables''. \emph{MathOverflow}, answer
478525, 9 September 2024. \url{https://mathoverflow.net/a/478525}.
\bibitem{Melbourne}
J. Melbourne and C. Roberto (2022). Quantitative form of Ball's cube slicing
in $\R^n$ and equality cases in the min-entropy power inequality.
\emph{Proceedings of the American Mathematical Society}, 150(8), 3595--3611.
\href{https://doi.org/10.1090/proc/15944}{doi:10.1090/proc/15944}.
\bibitem{Szarek}
S. J. Szarek (1976). On the best constants in the Khinchin inequality.
\emph{Studia Mathematica}, 58(2), 197--208.
\href{https://doi.org/10.4064/sm-58-2-197-208}{doi:10.4064/sm-58-2-197-208}.
\bibitem{YuLevel}
L. Yu (2026). On average distance, level-1 Fourier weight, and Chang's
lemma. \emph{arXiv}:2504.02593v2 [math.CO], 25 September 2026.
\url{https://arxiv.org/abs/2504.02593v2}.
\end{thebibliography}
\end{document}
